\chapter{Conclusion}
\label{ch:conclusion}

The conclusion is divided into three sections and summarizes the core findings of this thesis.
In \autoref{sec:platform-performance}, we assess the performance of the \tech{ln-history} platform,
specifically how well the platform architecture met its requirements in a production environment and how to improve it.

After that we delve into a discussion about the meaning of the analysis.
\autoref{sec:ln-gossip-protocol-conclusion} presents a conclusion about the \autoref{sec:metadata-analysis}.
Likewise, the \autoref{sec:ln-evolution} concludes the \autoref{sec:historic-gossip-analysis}.


\section{\tech{ln-history} Platform Performance}
\label{sec:platform-performance} 

The evaluation of the platform's functional and non-functional requirements demonstrates
that the \tech{ln-history} architecture is capable of supporting advanced, empirical network research.

Through its custom data model, the platform manages long-term data storage and persistence (\ref{req:persistance}), 
alongside fulfilling the dual-format constraint for preserving both raw wire bytes and parsed data (\ref{req:dual-format}). 

For historical analysis, the system achieves snapshot retrieval (\ref{req:snapshot}) times averaging under \qty{3}{\second}.
However, we observed that while database retrieval is exceptionally fast, 
the resulting payload for a modern network snapshot still requires approximately \qty{50}{\mega\byte}. 
Consequently, network bandwidth and not the database performance remains 
the primary temporal bottleneck for end-users executing these queries.

In terms of real-time operations, the platform demonstrated reliability, 
running stably without major outages throughout the 48-day collection window. 
The real-time ingestion pipeline successfully processed multiple gigabytes of gossip messages (\ref{req:ingestion}), 
consistently maintaining an ingestion lag of under \qty{5}{\minute} under normal network conditions
without blocking parallel analytical queries (\ref{req:concurrency}). 

Meeting the non-functional requirements self-hostability (\ref{req:self-host}) and 
long-term maintainability (\ref{req:maintainability}) proved more challenging. 
While the system can operate on small-scale hardware, processing heavy \gls{olap} queries and 
bulk ingestions pushed the server to its thermal limits.

Maintaining optimal query performance requires active administration.
The commands \code{ANALYZE VACUUM FULL} and \code{CLUSTER} are necessary to prevent index degradation over time. 
Those challenges lead to the decision to migrate the database to a dedicated datacenter \gls{vps} 
ensuring the reliability required for the empirical analysis in this thesis.


\section{Lightning Network Gossip Protocol}
\label{sec:ln-gossip-protocol-conclusion}


\subsection{Network Gossip Volume and Lag}

As \autoref{sec:metadata-analysis} shows, Lightning nodes do not inherently have a 
correctly synchronized or complete view of the network topology. 
Beyond that, the protocol has disparities in propagation latency depending on the message type. 
During data cleaning, more than half of the collected \code{node\_announcements} 
were broadcasted redundantly without containing any novel topological information. 
Given the propagation blind spots, the inherent latency variations, and the overhead of redundant metadata, 
the empirical analysis shows that researching and implementing optimizations to the gossip protocol is necessary for long-term scalability. 
Various approaches to address these architectural bottlenecks will be discussed in the upcoming \autoref{ch:future-work}.


\subsection{Peer Session Analysis}

This thesis contributes empirical observations regarding the peer management behaviour of default \gls{cln} nodes. 
While the \gls{ln} operates without an enforcement to forward gossip messages, 
active participation in the gossip protocol fundamentally relies on stable \gls{p2p} connections. 
The analysis demonstrated that peer session durations have an immense variance, ranging from minutes to multiple days. 
When coupled with the synchronization inaccuracies discussed previously, 
the connection instability can yield to missing gossip messages.

Nodes that offer routing services are economically incentivized to maintain highly reliable peer connections,
a dynamic empirically supported by our finding that total node capacity positively correlates with the session duration. 
By ensuring uninterrupted access to gossip messages, specifically \code{channel\_updates},
these enterprise-level hubs can minimize payment routing failures and optimize their overall service reliability.
Smaller nodes that have an unstable connection are more prone to payment failures. 
By missing critical \code{channel\_updates} during periods of downtime, 
these nodes risk to use outdated fee policies and inaccurate topological data during their pathfinding calculations.



\section{Lightning Network Evolution}
\label{sec:ln-evolution}

The evolution of the \gls{ln} is a complex topic and can be viewed from different angles.
As we will see, the network has grown and matured in the past eight years.
Nonetheless, there are still many open questions and possible improvements to be made.

Following the structure of the \autoref{sec:historic-gossip-analysis}, 
we first discuss the quality of the gossip dataset currently persisted in the \tech{ln-history} platform.
After that we will examine empirical measurements of all observed closed channels.
Lastly, we will dive into discussions about the centrality of the \gls{ln}.


\subsection{Collected Gossip Messages}

As mentioned in \autoref{sec:historic-gossip-analysis}, evaluating the complete 
repository of over \num{91.3} million historical gossip messages 
reveals two distinct phases of intensive data collection, 
separated by a multi-year observational blind spot between 2022 and mid-2025. 
Despite this gap in ephemeral \code{channel\_updates}, 
the structural integrity of the historical network graph can be approximated reconstructed 
because the foundational \code{channel\_announcements} were largely preserved. 
Furthermore, the \autoref{fig:gossip-messages-distribution} shows a hierarchy in network traffic composition. 
Dynamic \code{channel\_updates} dominate the broadcast volume, followed by informational \code{node\_announcements}, 
while structural \code{channel\_announcements} represent only a tiny fraction of the gossip protocol.


\subsection{Channels}

Payment channels form the core of the \gls{ln} protocol. 
In this section, we discuss the empirical findings of over \num{283226} closed channels.
The theoretical objective of a payment channel is to remain open indefinitely, 
allowing nodes to route an unlimited number of off-chain transactions while 
amortizing the initial on-chain setup fees. 
However, as \autoref{fig:channel-lifetime-distribution} illustrates, 
the empirical reality is different: of all channels that have ever been funded, 
more than \qty{75}{\percent} have been closed.

It is reasonable to assume that larger channels might have a longer lifespan than smaller ones. 
Node operators committing significant capital to open large channels 
likely plan their network topology more carefully. 
To test this hypothesis, we measured the correlation between channel capacity (\code{capacity\_sat}) 
and channel lifetime (\autoref{fig:channel-lifetime-correlation}).
Surprisingly, the data reveals no correlation between these two variables. 
A channel's financial size does not predict how long it will remain open.

Another interesting observation shows the scatter plot. 
Channel capacities with round numbers (e.g., \code{capacity\_sat} ending in $00000$) appear more often (see: purple vertical lines). 
This clustering indicates that funding amounts are rarely mathematically optimized for exact routing needs. 
Instead, they are often chosen based on human intuition or simple convenience.

Ultimately, channel management proves to be far more complex than 
simply opening a connection and leaving it alone. 
The lack of correlation to capacity shows that channels frequently require active adjustments, closures, and rebalancing. 
This reality affects all participants, regardless of whether they are a professional hub or an individual hobbyist.


\subsection{Centrality}

The academic discourse surrounding centrality in the \gls{ln} is extensive, 
compared to other research areas of the \gls{ln} protocol. 
Numerous studies have applied a wide range of graph-theoretic metrics 
to historic snapshots to quantify the network's topology (e.g., \cite{atmanaviciute_quantitative_2025}). 
The empirical data presented in \autoref{sec:centrality} aligns with this literature, 
confirming that the network is highly centralized according to \gls{bc}.

This thesis argues that the academic narrative is often one-sided. 
Applying centrality metrics to static snapshots to condemn the network for deviating from Bitcoin's decentralized ethos,
fails to critically evaluate the resulting operational value that dominant nodes provide.
From an economic and operational perspective, a centralized topology is a logical and arguably beneficial development. 
Professional routing hubs provide liquidity and continuous uptime necessary for a functional payment network. 
The presence of these dedicated backbone nodes is likely a driver behind the 
high payment success rates and overall service reliability that the \gls{ln} currently provides~\cite{random__reran_2024}.
Enforcing strict, uniform decentralization would likely increase routing failures.

Nonetheless, from the perspectives of privacy, censorship resistance, and systemic independence, 
this centralization of routing power definitely remains concerning. 
Relying on a handful of dominant hub operators contradicts the \gls{p2p} ethos of Bitcoin. 
These central hubs represent a significant vulnerability regarding user privacy. 
While Lightning payments utilize onion routing, central hubs process such a massive share of the total traffic 
that they are uniquely positioned to perform traffic analysis. 
By monitoring the flow, timing, and channel balances of adjacent packets, 
dominant routing nodes could theoretically try to trace and deanonymize payment flows, 
severely degrading the privacy guarantees originally engineered into the protocol.

Ultimately, the centrality of the \gls{ln} represents compromise that distributed systems have to face by design. 
The network has organically gravitated toward a topology that sacrifices a degree of ideological decentralization and 
strict privacy to achieve the operational efficiency, liquidity, and reliability required for real-world adoption.

\iffalse
\subsection{Functional Requirements}

\subsubsection{Data Storage \& Persistence (RF 1 \& 2)}
Thanks to the database centric layout with its custom data model,
the first two functional requirements RF1 and RF2 are met.

\subsubsection{Querying \& Historic Analysis (RF 3 - 5)}
To make \gls{olap} queries a lot of index maintanance is necessary.
The storage used by indexes reaches about the storage of the actual data
but is necessary to provide in depth analysis.

\paragraph{Snapshot Retrieval RF3}
The thrid functional requirement (RF3) stating that the plaform is able to retrive 
historic snapshots has been met very well. 
The pure retrieval of gossip messages from the database 
taks on average less than $\qty{3}{s}$. 

As the database utilizses its \gls{ram} to keep data that is often requested \textit{warm},
snapshot queries for similar timestamp can yield results even quicker.
As the analysis has shown snapshots with thousands of nodes and edges take up space.
A snapshot of the current lightning network from the \tech{ln-history} plaform,
which already reduces the number of gossip messages drastically and sends them in raw bytes,
still needs aroude $\qty{50}{\mega B}$, which depending on the network bandwidth
takes significantly more time to download than the actual retrieval from storage.

This main feature is working \textit{lightning} fast and shows that the \tech{ln-history}
platform is a significant upgrade for researchers when analysing historic snapshots.

\paragraph{Time-Window Retrieval RF4}
The fourth functional requirement (RF4) initially sounds interesting and close by
the snapshot retrieval but at the time of writing this thesis there has not been
many analysis ideas been found. 
The queries are, depending on the length of the time window,
rapid fast but still there is no use case yet.

\paragraph{Entity Filtering RF5}
This requirement is fulfilled and used in multiple queries specifically analysis
over channels or nodes.


\subsubsection{Real-Time Operations RF6 - RF10}
As shown in the first part of the analysis \autoref{sec:metadata-analysis},
the \tech{ln-history} platform has been running stable between  
%START_DATETIME_STR = "2026-01-14T00:00:00Z"
%END_DATETIME_STR = "2026-03-03T00:00:00Z"
.
In that time it processed multiple $\qty{}{\giga B}$ of gossip messages and
was able to ingest them.


\paragraph{Low-Latency Ingestion RF6}
RF6 which requires the database to reflect the live state of the network
by quickly inserting new gossip is reached.
The database thanks to batching gossip messages into a database transactions,
the ingestion throughput is not a problem and the \service{ln-history-database}
is capable to ingestion an order of magnitude more messages per second as normally
occur in the \gls{ln}

\paragraph{Concurrency RF7}
This requirement is almost fully met. In phases of exceptionally high workloads
onto the database (e.g. generation of many snapshots or complex \gls{olap} queries)
the insertion almost stops and an ingestion lag builds up.
But under normal conditions, the lag remains below $\qty{10}{s}$.

\paragraph{Multiple nodes RF8}
As of writing this thesis only two nodes have been collecting gossip messages.
The platform was capable of managing those, it should be possible to scale to at least
five concurrent running nodes.

\paragraph{Metadata RF9}
Thanks to the introduction of the \code{gossip\_inventory} and \code{gossip\_observations}
every collected gossip message immediatly gets an unique identifier \code{gossip\_id}.
Together with the \code{collector\_node\_id}, makes every gossip message unique.
Since every recieved gossip message gets its timestamp appended to it, it is really
easy to do propergation lag analysis over the gossip dataset. 


\paragraph{Batch ingestion RF10}
This script is doing its job but needs refactoring, since its generally a manaul 
process occuring rarely we consider this circumstance as fine.
Nonetheless, batch ingestions require the \service{gossip-processor} to stop
otherwise the ingestion will (depending on the size) take multiple days.


\subsection{Non-Functional Requirements}
The Non-Functional Requirements were defined by the following paragraphs and 
have sometimes been more challenging to meet than the functional ones.

\paragraph{Transparency \& Auditability (Open Source) NF 1}
This requirement forces the \tech{ln-history} platform to be fully open source
making it possible for other researchers to verify the results.
It met, by the fact that the source code of each component is openly available
\footnote{See https://github.com/ln-history}.


\paragraph{Deployability (Containerization) NF 2}
Thanks to the usage of Docker every component is containerized and working as intended.

\paragraph{Efficiency (Self-Hostability) NF 3}
This requirement was by far the most difficult to meet. 
We own a small scale server with a five-threaded CPU and DDR5 \gls{ram} as well
as modern NVMe SSD storage.
Nonetheless, during high \gls{olap} queries or at startup, where large amount
of missing gossip messages were being analysed / ingested, the server got heated up.
Since the platform ran without major outtage for 48 consecutive days, we consider this
requirement met, although for doing the analysis in chapter \autoref{ch:empirical-data-analysis},
we moved the database to a VPS in a datacenter, which in comparision proved so much more reliability.


\paragraph{Maintainability (long-term) NF 4}
Although the architecture strongly follows the \textit{Seperation of Concerns} pattern,
having multiple (independent) services working together, the maintainability can still be challenging.
Different gossip formats require delicate handling as updating already inserted gossip messages can get really difficult.
Also changing the interface of one service can have devistating cascading effects,
where the other components in the processing pipeline crash.
During development those constaintraints often lead to hours or days of troubleshooting.

Also the \service{ln-history-database} needs regular active maintainance to provide the expected service.
Especially the indixes and what postgres thinks of them deviates over time.
Using \code{ANALYZE VACUUM FULL} is a must, as well as regular \code{CLUSTER} commands  
\fi